\documentclass[10pt,a4paper]{amsart}
\usepackage[margin=19mm]{geometry}
\usepackage{amsmath,amssymb,mathtools}
\usepackage[scr=rsfs,cal=euler]{mathalfa}
\usepackage{xfrac}
\usepackage[unicode,hidelinks,bookmarks=false]{hyperref}
\newcommand{\ie}{i.e.\ }
\newcommand{\cf}{cf.\ }
\newcommand{\resp}{respectively\ }
\newcommand{\Subsection}{\subsection}
\providecommand{\nobreakdash}{\nobreak\hbox{-}}
\setlength{\emergencystretch}{2em}
\multlinegap0pt
\usepackage{bm}
\usepackage{bbm}
\usepackage{dsfont}
\usepackage{xspace}
\usepackage{cancel}
\usepackage{mathtools}
\usepackage{amsthm}
\makeatletter
\@ifundefined{lemma}{\newtheorem{lemma}{Lemma}}{}
\makeatother
\newcommand{\ceil}[1]{\left\lceil #1 \right\rceil}
\title[An algorithm for the Product-Sum Equality]{An algorithm for the Product-Sum Equality}
\author{Hlib Husarov, Eberhard Mayerhofer}
\date{Source: arXiv:2508.09647v1 (CC BY 4.0)}
\begin{document}
\maketitle

\noindent\emph{Source and licence.} An algorithm for the Product-Sum Equality. Version arXiv:2508.09647v1 (CC BY 4.0), distributed under the Creative Commons Attribution 4.0 International licence, as recorded in the arXiv licence metadata for this exact version and shown on the abstract page https://arxiv.org/abs/2508.09647v1. Full rights notice: licenses/arxiv-2508.09647-CC-BY-4.0.txt.\par\smallskip

\noindent\textit{Editorial scope.} Counted unit: the Lemma whose source label is \texttt{lem: log length} in the article (source0.tex lines 479--488, source label \texttt{lem: log length}), delivered together with the article's own complete original proof of it (source0.tex lines 492--524, one \texttt{proof} environment of 33 lines). No mathematics is written, repaired or re-derived: statement and proof are the article's own text line for line. The proof is detached from the statement in the source: the article states the result at lines 479--488 and prints its proof at lines 492--524, and the proof environment's own header names the statement, so the association is the article's own. Required context retained uncounted: none. Every definition, display and label of those lines is retained unchanged. Omitted from this package and not redistributed: 1--478, 489--491, 525--602. Nothing inside a retained excerpt is omitted or altered: every retained body is delivered exactly as the article's lines are written, so no omission receipt of any kind accompanies this package. Citations: the retained excerpts carry no citation command at all, so no bibliography entry of the article is transcribed and no reference is redistributed. Reference adaptation: none was needed and none was made; every reference inside the delivered unit resolves inside it, so no unresolved reference or citation is produced. Printed slips kept literal: no printed word, symbol, hypothesis or number of the counted statement or of its proof was altered. Layout and class: the article's own class is \texttt{article} with packages including amssymb, latexsym, amsmath, amsthm, amsfonts, epsfig, mathtools, xcolor, listings, url, natbib; the wrapper is the lane's pinned \texttt{amsart} editorial wrapper at 10pt on A4, which restates the theorem-like environments the retained text needs and adds the article's own macro definitions that the retained text needs (\texttt{\textbackslash ceil}), with extra wrapper packages (bm, bbm, dsfont, xspace, cancel, mathtools). The delivered numbering is the unit's own. Assets: no retained span contains a figure environment, a graphics-inclusion command, a PDF-page inclusion, a table or a TikZ picture; no image file is copied into this package, the package declares no asset dependency and no image file is redistributed with it. Licence: version arXiv:2508.09647v1 of this article is distributed under the Creative Commons Attribution 4.0 International licence, as recorded in the arXiv licence metadata for this exact version and shown on the abstract page https://arxiv.org/abs/2508.09647v1. A full rights notice is delivered at licenses/arxiv-2508.09647-CC-BY-4.0.txt. Characters: every retained excerpt is pure ASCII, so the delivered engine, XeLaTeX with its default Latin Modern text font, reports no missing character.
\begin{lemma}\label{lem: log length}
Let $n\geq 3$ and let $a$ satisfy the inequality
\begin{equation}\label{ineq: sumprod}
\prod_{i=1}^n a_i\leq\sum_{i=1}^n a_i.
\end{equation}
For $2\leq m\leq n$, let $k_{\star,m}$ be the greatest integer $k$ for which $a_k=m$. Then
\begin{equation}\label{eq: log}
k_{\star,m}< \log_m(n)+1.
\end{equation}
\end{lemma}

\begin{proof}[Proof of Lemma \ref{lem: log length}.]
First, we claim that for $n\geq 3$, $k_{\star,2}<n$ (and thus $k_{\star,m}<n$ for any $m>1$). In fact, this can be directly checked to hold for $3\leq n\leq 4$. Furthermore, for $n\geq 5$, we know by induction that $2^n>n^2$. Assume, for a contradiction, $k_{\star,2}=n$. Then $2^n\leq \prod_{i=1}^n a_i\leq\sum_{i=1}^na_i\leq n^2$, which is impossible. 

Thus we have to prove that if the product of $n \geq 3$ natural numbers $a_1 \geq a_2 \geq \cdots \geq a_n \geq 1$ satisfies \eqref{ineq: sumprod} and precisely $1\leq k<n$ of them are greater than or equal to $m$:
    \begin{gather}
        a_1 \geq a_2 \geq \cdots \geq a_k \geq m \label{geq_m} \\
        m > a_{k+1} \geq a_{k+2} \geq \cdots \geq a_n \label{less_than_m}
    \end{gather}
    Then:
    \[ k \leq \ceil {\log_m n}. \]

Proof: Since $a_1 \geq a_2 \geq \cdots \geq a_k$, we have:
    \begin{equation}\label{ineq: sumprodxx} 
        a_1 \ a_k^{k-1} \leq \prod_{i=1}^k a_i \leq \prod_{i=1}^n a_i\leq  \sum_{i=1}^n a_i
    \end{equation}
    From \eqref{geq_m} and \eqref{less_than_m}, we have:
    \begin{equation}\label{not_strict}
        \sum_{i=1}^n a_i < a_1 \ k +  m \ (n-k).
    \end{equation}
    Substituting this into \eqref{ineq: sumprodxx} we get:
    \[
        a_1 \ a_k^{k-1} < a_1 \ k + m \ (n-k).
    \]
    Dividing both sides by $a_1$, we get:
    \[
        a_k^{k-1} < k + \frac{m}{a_1} \ (n - k)\leq n,
    \]
the last inequality holding due to $a_1 \geq m$. Thus,
    \[
        k-1 < \log_{a_k} n
    \]
    From \eqref{geq_m}, we know that $a_k \geq m$, thus    $k < \log_m n + 1$.
\end{proof}

\end{document}
